\chapter{Introdução}
\label{sec:introducao}

% Fórmula de Head, via Bellemare: gancho, pergunta, antecedentes,
% contribuição, roteiro. Os elementos que a NBR 15287 pede a um projeto
% (objetivos, hipóteses, justificativa) vêm em seções próprias, depois.

Em fevereiro de 2009, uma empresa de aviação agrícola foi contratada para
pulverizar mil hectares de bananais na Chapada do Apodi, no Ceará. A primeira
aplicação daquele ano tomou cinco dias e vinte e nove voos, com o avião
decolando pela manhã e no fim da tarde do campo de pouso da Chapada
\cite[p.~171--173]{marinho2011dimensao}. Moradores descreviam galinhas mortas no
quintal e roupas no varal com cheiro de veneno.\footnote{Relato da pesquisadora
Raquel Rigotto e de uma moradora da comunidade do Tomé à Agência Brasil, 16 de
julho de 2016.} No ano seguinte, o líder comunitário que fotografava os voos,
José Maria Filho, o Zé Maria do Tomé, foi assassinado. Em janeiro de 2019, a lei
estadual que leva o seu nome proibiu a pulverização aérea de agrotóxicos em
todo o Ceará, e o estado se tornou o primeiro do país a fazê-lo.

A proibição é uma política pública com objetivo declarado de saúde, e não foi
avaliada. Sabe-se que foi aprovada, que o Supremo Tribunal Federal a julgou
constitucional e que, quase seis anos depois, o legislador estadual abriu uma
exceção para drones. Não se sabe se melhorou a saúde de quem vivia ao lado dos
bananais, nem o que custou à produção. Os dois lados do debate oferecem números
sem grupo de comparação: produtores falam em queda de produção, defensores da
lei mostram que a produção estadual de banana subiu.

\section{Pergunta de pesquisa}

A proibição da pulverização aérea de agrotóxicos no Ceará melhorou a saúde ao
nascer nos municípios da Chapada do Apodi, e a que custo para a produção
agrícola?

\section{Antecedentes}

Três conjuntos de trabalhos estão mais próximos desta pergunta.

O primeiro documenta o dano na própria Chapada, sem identificação causal.
Resíduos de fungicidas foram medidos em águas subterrâneas
\cite{tramas2016parecer}. \citeonline{rigotto2013trends} encontram taxas de
internação e de morte por neoplasia mais altas nos municípios da fruticultura
do que em municípios de agricultura familiar. \citeonline{aguiar2017} descreve
casos de malformação congênita em uma comunidade de dois mil e quinhentos
habitantes.

O segundo é a economia brasileira dos agrotóxicos. \citeonline{dias2023down}
mostram, com identificação crível, que a expansão do glifosato na soja elevou a
mortalidade infantil a jusante das áreas de uso. \citeonline{rodrigues2022tese}
e \citeonline{januzzi2025} tentam medir efeitos em escala nacional com proxies
municipais de uso; essas proxies não distinguem o método de aplicação e não
enxergam a fruticultura irrigada, e a segunda autora não encontra efeito sobre
nascimentos.

O terceiro é o análogo internacional. \citeonline{calzada2023bananas} estudam a
fumigação aérea de bananais no Equador e encontram déficit de peso ao nascer
entre filhos de mães que moram a até 150 metros das plantações.
\citeonline{camacho2017health} encontram aumento de perdas gestacionais com a
aspersão aérea de herbicida na Colômbia.

Nenhum desses trabalhos avalia uma política que tenha retirado a aplicação
aérea.

\section{Contribuição}

O projeto contribui em três frentes.

\begin{enumerate}
\item \textbf{Identificação.} A variação vem de uma lei com data e fronteira, e
não de diferenças de uso inferidas de área plantada. A Chapada do Apodi
atravessa a divisa com o Rio Grande do Norte, que não proibiu, o que fornece um
grupo de comparação com o mesmo clima e o mesmo aquífero.
\item \textbf{Objeto e dados.} Avalia-se a proibição de um \emph{método} de
aplicação, em fruticultura irrigada, com fungicidas: um caso que a literatura
brasileira, centrada em herbicida e grãos, não cobre. E o projeto reúne, por
pedidos de acesso à informação, os registros administrativos da aplicação aérea
nos anos anteriores à lei.
\item \textbf{Mecanismo.} A teoria da mudança da lei é testada elo a elo: se os
voos cessaram, o que os substituiu, o que aconteceu com as intoxicações de
trabalhadores, com a produção e com a perda fetal. O efeito sobre o nascimento
só é atribuído à redução da deriva depois de examinados os caminhos
concorrentes.
\end{enumerate}

O projeto também delimita o que não promete. A exposição era localizada, em
torno de mil hectares, e os dados de nascimento são municipais. O cálculo de
poder do Capítulo~\ref{sec:poder} mostra que o efeito esperado sobre o peso
médio municipal é de um a dois gramas no caso central, muito abaixo do que o
desenho principal detecta. O plano responde a esse limite com um desenho de
calendário, com desfechos de perda fetal e com a regra, fixada antes de
estimar, de que o resultado é o intervalo de confiança e o que ele exclui.

\section{Objetivos}

\subsection{Objetivo geral}

Avaliar o impacto da proibição da pulverização aérea de agrotóxicos no Ceará
(Lei nº~16.820/2019) sobre a saúde ao nascer nos municípios da Chapada do
Apodi.

\subsection{Objetivos específicos}

\begin{enumerate}[label=\alph*)]
\item caracterizar a política e reconstruir a sua teoria da mudança;
\item medir, com registros administrativos, a extensão da aplicação aérea nos
anos anteriores à lei;
\item estimar o efeito da lei sobre o peso ao nascer, a duração da gestação e a
perda fetal, comparando os dois lados da Chapada;
\item verificar os caminhos concorrentes: substituição de tecnologia de
aplicação, intoxicações de trabalhadores, produção e emprego;
\item examinar se o efeito difere entre municípios de agricultura familiar e não
familiar;
\item comparar benefícios e custos da proibição e discutir a escolha do
instrumento.
\end{enumerate}

\section{Hipóteses}

\begin{description}
\item[H1] A proibição melhorou os desfechos de nascimento nos municípios
cearenses da Chapada, em relação aos potiguares.
\item[H2] O efeito, se existir, concentra-se nas gestações que atravessam a
estação de aplicação e nos municípios em que a agricultura não familiar aplicava
por aeronave.
\item[H3] A substituição do avião pela aplicação terrestre não reduziu, e pode
ter elevado, as intoxicações de trabalhadores.
\item[H4] O custo da proibição em área e produção de banana foi pequeno.
\end{description}

As hipóteses H1 e H2 só se testam com sentido se a aplicação aérea estava em
curso quando a lei entrou em vigor. Essa premissa é objeto de medida, e não de
suposição (Seção~\ref{sec:quando}).

\section{Justificativa}

A avaliação se justifica pelos critérios usuais para decidir quando avaliar uma
política \cite{gertler2016}. A política é inovadora: o Ceará foi o primeiro
estado a proibir. É replicável: o Supremo Tribunal Federal reconheceu a
competência estadual para legislar sobre a matéria. Não foi testada: os números
em circulação não têm grupo de comparação. E o próprio legislador a reviu em
2024, ao admitir drones, sem evidência sobre o que os cinco anos anteriores
produziram.

Há também uma razão de método. A evidência brasileira sobre agrotóxicos e
saúde apoia-se em medidas indiretas de uso. Uma lei que retira um método de
aplicação, numa região onde o dano foi descrito por quinze anos, com um vizinho
que não a adotou, é uma oportunidade rara de observar o efeito de uma mudança
datada.

O projeto é viável. Os dados de nascimento, mortalidade, intoxicação e produção
são públicos, e os registros da aplicação aérea são solicitáveis pela Lei de
Acesso à Informação.

\section{Organização do trabalho}

O Capítulo~\ref{sec:politica} apresenta o contexto institucional: a lei, a
Chapada e a prática proibida. O Capítulo~\ref{sec:literatura} revê a
literatura. O Capítulo~\ref{sec:desenho} reconstrói a teoria da mudança da
política e avalia o seu desenho. O Capítulo~\ref{sec:avaliacao} traz a
estratégia empírica: parâmetro, identificação, estimador, amostra e dados. O
Capítulo~\ref{sec:poder} traz o cálculo de poder. O Capítulo~\ref{sec:cb}
esboça o segundo ensaio, de custo-benefício e escolha de instrumento. O
Capítulo~\ref{sec:cronograma} apresenta o cronograma e os recursos, e o
Capítulo~\ref{sec:conclusao}, as considerações finais.
